\section{Discussion}

Across the six primary families, the common pattern is conditional rather than universal. Typing/scoping matters when it changes which uncertainty, receipt, certificate, or experiment is relevant to the decision. The matched-information ablations make that interpretation stronger than a comparison in which the typed arm simply sees richer state.

Equally important is the no-value behavior. N4-F3's exact tie shows that the mechanism does not need to win everywhere to support its intended claim; in the regime where remint/transport has no informational value, it does nothing and ties. N1-C pushes the same principle further: typed state can matter while the proposed policy adds nothing beyond an ideal donor. The method claim should stop at the state representation in that case.

Hostile controls are useful because they test whether a synthetic world actually exercises the proposed failure mode. They do not make the world realistic. A model can perform well on all six constructions and still fail on real research interfaces for reasons absent from the generators: mis-specified state types, noisy provenance, nonstationary costs, strategic adversaries, or poor semantic extraction from raw evidence.

The next scientifically discriminating step is therefore transfer, not more synthetic variants of the same traps. A real-domain instantiation would need independently auditable state labels and matched-information baselines rather than silently substituting a richer system for the typed mechanism.